% ===== PRÉAMBULE CANONIQUE — à coller VERBATIM en tête de CHAQUE .tex =====
\documentclass[11pt,a4paper]{article}
\usepackage[utf8]{inputenc}\usepackage[T1]{fontenc}\usepackage{lmodern}
\usepackage[french]{babel}
\usepackage{amsmath,amssymb,mathtools}\usepackage{icomma}
\usepackage{siunitx}\sisetup{locale=FR}  % \qty{}{}, \si{}, \ang{}
\usepackage{xcolor}\usepackage{colortbl}
\usepackage{tikz}\usepackage{pgfplots}\pgfplotsset{compat=1.18}
\usepackage[siunitx,europeanresistors]{circuitikz} % circuits (élec.)
\usepackage[most]{tcolorbox}\usepackage{enumitem}\usepackage{array,booktabs}
\usepackage[a4paper,margin=2.2cm]{geometry}
\usetikzlibrary{arrows.meta,calc,angles,quotes,positioning,patterns,%
  backgrounds,shapes.geometric,decorations.pathmorphing,decorations.markings,intersections}
\definecolor{primary}{RGB}{222,49,99}\definecolor{primarydark}{RGB}{150,22,55}
\definecolor{defcol}{RGB}{0,114,178}\definecolor{methcol}{RGB}{52,92,148}
\definecolor{excol}{RGB}{0,135,90}\definecolor{attncol}{RGB}{200,40,55}
\tcbset{boxbase/.style={enhanced,breakable,boxrule=0.9pt,arc=2.5pt,
  left=3mm,right=3mm,top=2.3mm,bottom=2.3mm,fonttitle=\bfseries,coltitle=white}}
% --- boîtes TUTORIEL (noms EXACTS, ne pas renommer) ---
\newtcolorbox{cle}[1][]{boxbase,colback=defcol!6,colframe=defcol,
  title={\textsc{À retenir}\ifx\relax#1\relax\relax\else\textmd{ : \itshape #1}\fi}}
\newtcolorbox{methode}[1][]{boxbase,colback=methcol!7,colframe=methcol,sharp corners,
  title={\textsc{Méthode}\ifx\relax#1\relax\relax\else\textmd{ --- \itshape #1}\fi}}
\newtcolorbox{exemplebox}[1][]{boxbase,colback=excol!5,colframe=excol,
  title={\textsc{Exemple}\ifx\relax#1\relax\relax\else\textmd{ : \itshape #1}\fi}}
\newtcolorbox{piege}[1][]{boxbase,colback=attncol!6,colframe=attncol,
  title={\textsc{Piège}\ifx\relax#1\relax\relax\else\textmd{ : \itshape #1}\fi}}
\newtcolorbox{remarque}[1][]{boxbase,colback=black!4,colframe=black!45,
  title={\textsc{Remarque}\ifx\relax#1\relax\relax\else\textmd{ : \itshape #1}\fi}}
% --- boîtes EXERCICE / CORRIGÉ (noms EXACTS, auto-comptées) ---
\newtcolorbox[auto counter]{exo}[1][]{boxbase,colback=primary!4,colframe=primary,
  title={\textsc{Exercice}~\thetcbcounter\ifx\relax#1\relax\relax\else\textmd{ --- \itshape #1}\fi}}
\newtcolorbox[auto counter]{corrige}[1][]{boxbase,colback=excol!4,colframe=excol,
  title={\textsc{Corrigé}~\thetcbcounter\ifx\relax#1\relax\relax\else\textmd{ --- \itshape #1}\fi}}
\newcommand{\vv}[1]{\vec{#1}}\newcommand{\ds}{\displaystyle}
\newcommand{\R}{\mathbb{R}}\newcommand{\N}{\mathbb{N}}\newcommand{\Z}{\mathbb{Z}}
% (un \usepackage{fancyhdr}+en-tête est possible ; il est ignoré côté web)
% =========================================================================
\begin{document}
\begin{corrige}[poids : application directe]
On applique $G=m\,g$ avec $m$ déjà en kilogrammes.
\begin{enumerate}[label=\alph*)]
  \item $G = 3\times 10 = \qty{30}{N}$.
  \item $G = 3\times 9{,}81 = \qty{29.43}{N}\approx\qty{29.4}{N}$.
\end{enumerate}
Les deux valeurs sont proches : arrondir $g$ à \qty{10}{N/kg} donne une bonne estimation.
\end{corrige}

\begin{corrige}[convertir avant de calculer]
\begin{enumerate}[label=\alph*)]
  \item $m = \qty{250}{g} = \dfrac{250}{1000}\ \si{kg} = \qty{0.250}{kg}$.
  \item $G = m\,g = 0{,}250\times 10 = \qty{2.5}{N}$.
\end{enumerate}
Piège évité : sans conversion, on aurait écrit $250\times 10=\qty{2500}{N}$, soit 1000 fois trop.
\end{corrige}

\begin{corrige}[réaction d'une table]
\begin{enumerate}[label=\alph*)]
  \item Le \textbf{poids} $\vv G$ (vertical, vers le bas) et la \textbf{réaction normale} $\vv R$ de la
        table (verticale, vers le haut).
  \item Équilibre de translation : $\vv G + \vv R = \vv 0 \Rightarrow \vv R = -\vv G$. En intensité,
        $R = G = \qty{15}{N}$, dirigée vers le haut.
\end{enumerate}
\end{corrige}

\begin{corrige}[objet suspendu à un fil]
\begin{enumerate}[label=\alph*)]
  \item $G = m\,g = 2\times 10 = \qty{20}{N}$.
  \item L'objet est en équilibre sous deux forces : $\vv T$ (vers le haut) et $\vv G$ (vers le bas).
        Donc $T = G = \qty{20}{N}$. Le dynamomètre du fil afficherait \qty{20}{N}.
\end{enumerate}
\end{corrige}

\begin{corrige}[le même corps, deux astres]
\begin{enumerate}[label=\alph*)]
  \item Sur la Terre : $G_{\text{T}} = 5\times 10 = \qty{50}{N}$. Sur la Lune :
        $G_{\text{L}} = 5\times 1{,}6 = \qty{8}{N}$.
  \item La masse $m=\qty{5}{kg}$ est une quantité de matière : elle est \textbf{identique} partout.
        Seul le poids change ; il est \emph{plus grand sur la Terre} (là où $g$ est plus grand).
\end{enumerate}
\end{corrige}

\end{document}
